\documentclass[aspectratio=169]{beamer}
\usetheme{Boadilla}
\usepackage{booktabs}

\title[Self-Healing Polymers]{Self-Healing Polymers for Flexible Electronics}
\subtitle{Department seminar series}
\author[T. Herrera]{Dr.\ Tomas Herrera}
\institute[Coastal Tech]{Department of Materials Science, Coastal Technical University}
\date{9 July 2026}

\begin{document}

\begin{frame}
  \titlepage
\end{frame}

\begin{frame}{Outline}
  \tableofcontents
\end{frame}

\section{Healing chemistry}

\begin{frame}{How the network heals itself}
  \begin{itemize}
    \item Dynamic disulfide bonds exchange at room temperature, so a cut
          interface re-forms crosslinks within hours.
    \item Hydrogen-bonding side chains provide fast initial tack while the
          covalent network recovers.
    \item No added catalyst or heat: healing proceeds under ambient conditions.
  \end{itemize}
\end{frame}

\section{Device performance}

\begin{frame}{Performance after damage}
  \begin{table}
    \centering
    \begin{tabular}{lrr}
      \toprule
      Property & Pristine & Healed (24 h) \\
      \midrule
      Tensile strength (MPa) & 6.8 & 6.1 \\
      Elongation at break & 640\% & 590\% \\
      Conductivity of trace (S/cm) & 2100 & 1870 \\
      \bottomrule
    \end{tabular}
  \end{table}
  \vspace{0.5em}
  Strain sensors built on the healed substrate survive 10,000 bending
  cycles with under 4\% drift.
\end{frame}

\begin{frame}{Closing}
  \begin{center}
    {\Large Thank you}\\[1em]
    Collaboration welcome: preprint and data at coastaltech.example/herrera-lab
  \end{center}
\end{frame}

\end{document}
